\section{Introduction}

Large Language Model (LLM) development is shifting from supervised pre-training toward post-training reinforcement learning (RL). Recent work in Reinforcement Learning has demonstrated that scaling RL compute together with test-time compute is a highly effective way to improve model intelligence \citep{deepseek2024r1, openaio1, cobbe2021training, lightman2023let}.
% The prevailing paradigm for training these models relies on synchronous, ``interleaved'' data generation and training. In this setting, the policy model generates a batch of responses, and optimization steps occur only after the entire batch is collected \citep{ouyang2022training, rafailov2024direct}. While the synchronous strategy has been widely adopted~\cite{}, this approach suffers severely from \textit{imbalanced generation overhead}. In agentic and coding tasks, generation lengths vary drastically; a short ``thought'' process might take seconds, while a complex reasoning chain requires minutes. Consequently, in a synchronous batch, the entire GPU cluster is forced to idle, waiting for the slowest ``straggler'' to finish before the gradient update can proceed \citep{shao2024deepseekv3, kwon2023efficient, yu2022orca}. This inefficiency motivates a shift toward Asynchronous Reinforcement Learning, where data is consumed and trained upon as soon as it is generated, thereby maximizing hardware utilization \citep{mnih2016asynchronous, liang2018rllib, hoffman2020acme}.
Most LLM RL pipelines remain synchronous and interleaved: the policy generates a batch of rollouts, and optimization starts only after the entire batch is collected \citep{ouyang2022training, rafailov2024direct}. 

For agentic and coding workloads, rollout lengths are highly variable, so short trajectories finish quickly while long ones become stragglers; as a result, large portions of the GPU cluster idle while waiting for the slowest rollouts \citep{shao2024deepseekv3, kwon2023efficient, yu2022orca}. Asynchronous RL mitigates this \textit{imbalanced generation overhead} by consuming rollouts continuously as they arrive, improving utilization and wall-clock efficiency \citep{mnih2016asynchronous, liang2018rllib, hoffman2020acme}.


% However, transitioning to asynchronous RL brings additional challenges. First, each response can be generated by multiple versions of the rollout model, which leads to more unpredictable and severe off-policy and hards training instability. 
% Second, the widely used Group Relative Policy Optimization (GRPO) \citep{deepseek2024math, wang2022self} has limitations in asynchronous settings. GRPO eliminates the need for a value function by normalizing rewards within a group of outputs generated from the same prompt. While effective in synchronous settings, GRPO is structurally antagonistic to asynchronous training for two primary reasons. 
% First, it introduces latency-induced off-policy behavior: waiting to collect a full group of diverse-length outputs forces the ``fastest'' generated responses to wait in a queue for the ``slowest'' response. By the time the group is complete, the learning policy $\pi_\theta$ may have drifted significantly, rendering the collected data heavily off-policy \citep{sutton2018reinforcement}.
% Second, it suffers from incompatibility with online and agentic tasks: in real-world online learning or multi-agent environments, the system typically receives a single trajectory feedback per prompt, making group-wise statistical baselines mathematically infeasible \citep{schulman2017proximal, yao2022react, nakano2021webgpt}.
However, asynchrony introduces two challenges. 
First, each trajectory can be generated by multiple versions of the old rollout model, which leads to more unpredictable and severe off-policy, and thus harms the training stability. Previous works~\citep{fu2025areal,noukhovitch2024asynchronous} make attempts for asynchronous RL but mainly focus on efficiency optimization rather than effectiveness.
% policy lag increases off-policy drift and can destabilize both policy and value learning. 
Second, group-wise methods such as GRPO \citep{deepseek2024math, wang2022self} are mismatched to asynchronous training. 
GRPO samples a group of responses for each prompt and uses the group-level average for advantage estimation.
The group-wise sampling induces latency-driven off-policy behavior because the group has to wait for the slower one to finish before fed into training. In addition, group-wise sampling is incompatible with online or complex agentic settings where the environment often provides only a single trajectory feedback per prompt \citep{sutton2018reinforcement, schulman2017proximal, yao2022react, nakano2021webgpt}.
% forming a group reintroduces waiting and staleness, and group-normalized baselines are incompatible with online/agentic settings where the environment often provides only a single trajectory feedback per prompt \citep{sutton2018reinforcement, schulman2017proximal, yao2022react, nakano2021webgpt}.




% In this paper, we aim to advance asynchronous RL by improving training stability and reducing the off-policy drift under policy lag. We propose \model,  a single-rollout asynchronous RL framework. Unlike group-based methods (e.g., GRPO), \model removes group-wise dependencies and updates the policy immediately after a single trajectory, maximizing throughput while minimizing lag-induced off-policy effects. 
% Abandoning the group-based baseline of GRPO necessitates a return to value-based critics, which are notoriously unstable in LLM training \citep{zheng2024secrets, engstrom2020implementation, andrychowicz2021matters}. To make single-rollout Async RL viable, we introduce a suite of stabilization techniques centered on training a robust Value Model:

In this paper, we propose Single-rollout Asynchronous Optimization (\model) for  agentic RL. 
It keeps asynchronous RL training  stable and effective under policy lag while preserving the efficiency of asynchrony.
%It replaces group-wise sampling such as GRPO with single-rollout updates. 
Instead of group-wise sampling, such as GRPO, \model uses single-rollout updates.
To make this setting practical, it also introduces effective value-model training strategies. Our contributions are as follows:

\begin{itemize}[leftmargin=*,itemsep=0pt,parsep=0.2em,topsep=0.3em,partopsep=0.3em]
\item To stabilize training under varied policy lag, we use token-level importance sampling strategy. It directly uses the log-probabilities from the rollout engine and applies stricter double-sided token-level clipping and masking.

\item To reduce off-policy effects, we use one single rollout sampling for each prompt instead of group-wise sampling previously populated by GRPO. 
To further make this setting practical in agentic RL, we improve the value model process. Specifically, we update the critic more frequent than the actor and fine-tune the value model with frozen attention.

\item To handle multi-turn agent trajectories with interleaved environment feedback, we derive a skip-observation token-level GAE estimator. 
It computes advantages across action-to-action boundaries. It also avoids propagating noise through observation tokens that are not generated by the model.

\end{itemize}


We evaluate \model{} on agentic coding and math reasoning benchmarks, including SWE-Bench Verified~\citep{jimenez2023swe}, AIME2025~\citep{balunovic2025matharena}, BeyondAIME~\citep{bytedance_seed_2025_beyondaime}, HMMT\citep{balunovic2025matharena}, and IMOAnswerBench\citep{luong-etal-2025-towards}.
The results demonstrate that our asynchronous RL design can stably train for around one thousand steps and achieves consistently better performance than improved GRPO. In addition, we show that the single-rollout strategy in \model is uniquely suited for simulated online learning, where it can adapt to dynamic environmental changes. 
